\documentclass[12pt, a4paper]{article}

% --- Encoding & Fonts ---
\usepackage[utf8]{inputenc}
\usepackage[T1]{fontenc}
\usepackage{lmodern}

% --- Mathematics ---
\usepackage{amsmath}
\usepackage{amssymb}
\usepackage{amsthm}
\usepackage{mathtools}

% --- Layout ---
\usepackage[margin=1in]{geometry}

% --- Lists ---
\usepackage{enumitem}

% --- Tables ---
\usepackage{booktabs}
\usepackage{array}

% --- Graphics ---
\usepackage{graphicx}

% --- Headers / Footers ---
\usepackage{fancyhdr}
\pagestyle{fancy}
\fancyhf{}
\fancyhead[L]{\leftmark}
\fancyhead[R]{\thepage}
\renewcommand{\headrulewidth}{0.4pt}

% --- Hyperlinks & Cross-references ---
\usepackage[colorlinks=true, linkcolor=blue, urlcolor=blue, citecolor=blue]{hyperref}

% --- Misc ---
\usepackage{xcolor}

% --- Custom commands (safety definitions) ---
\providecommand{\operatorname}[1]{\mathop{\mathrm{#1}}\nolimits}

\title{Learning Module: Sdos3}
\author{}
\date{\today}

\begin{document}

\maketitle
\thispagestyle{empty}
\newpage

\tableofcontents
\newpage

% ============================================================
% SECTION  1. Background and Need for Requirements
% ============================================================

% !!!  REVIEW FLAG: 1. Background and Need for Requirements  !!!
% [knowledge_assumption] Provide a brief preface within the theory section that introduces risk concepts or add a reference to a practice section on risk management.
% [undefined_term] Add a brief definition of error propagation (e.g., how errors in requirements can cascade and amplify costs) in the Formal Definition or an accompanying paragraph.
% !!! END REVIEW FLAGS !!!
\section{Background and Need for Requirements}
\subsection{Before We Begin}
A software development lifecycle is a series of stages that guide a product from an initial idea to a finished system. One early stage is requirements analysis, where we determine what the system must do. This stage sits after the initial business idea and before design and coding begin. Understanding this placement helps explain why a clear set of requirements is crucial. With that context, we can now explore the relationship between business goals and software requirements.

\subsection{Motivation}
Writing a recipe: the business goal is the desired dish, while the software requirements are the exact list of ingredients and stepbystep instructions needed to make it. Just as a chef needs a precise recipe to produce the intended meal, a development team needs wellcrafted requirements to deliver a system that meets the intended impact.

\subsection{Intuition}
Think of a business goal as the big picture outcome a company wants, such as ``increase sales'' or ``reduce support costs.'' To achieve that outcome, the software must perform specific actions, like recording transactions or generating reports. Requirements are the concrete, testable statements that describe those actions. By translating highlevel goals into clear requirements, everyone knows exactly what needs to be built and can verify that the final product delivers the desired result.

\subsection{Formal Definition}
The Institute of Electrical and Electronics Engineers (IEEE) defines a software requirement as a condition or capability that a system must possess. A business goal is an overarching objective that reflects the desired impact on the organization or society. A Software Requirements Specification (SRS) is a documented collection of all software requirements, created after the requirements analysis phase. The SRS serves three main purposes: it clarifies what users need (since users may not know all details themselves), it establishes a mutual agreement between the user and the supplier, and it provides a reference point for later validation of the delivered product. Because requirements are often hidden in large systems, errors in the early stages can become very expensive to fix later; a highquality SRS therefore helps keep development cost and effort low.

\subsection{Worked Example}
\textbf{Problem statement:} A retail store wants to increase its monthly sales by 10 %.

\begin{enumerate}
    \item \textbf{Identify the business goal.}\\
    The store's business goal is ``increase monthly sales by 10 %.''
    \item \textbf{Ask what the software must do to support this increase.}\\
    One way to support higher sales is to ensure every transaction is captured accurately, enabling analysis of purchasing patterns.
    \item \textbf{Formulate the first concrete requirement.}\\
    \emph{Requirement 1}: ``The system shall record each transaction, including item details, quantity, price, and timestamp.''
    \item \textbf{Identify another capability that helps monitor progress toward the goal.}\\
    The store needs regular insight into sales performance to act quickly on trends.
    \item \textbf{Formulate the second concrete requirement.}\\
    \emph{Requirement 2}: ``The system shall generate a daily sales report that summarizes total revenue, number of transactions, and bestselling items.''
\end{enumerate}

\textbf{Final answer:} The business goal of a 10 % sales increase is translated into two software requirements: (1) record each transaction, and (2) generate a daily sales report.

\subsection{Key Takeaway}
A business goal states the intended impact, whereas a software requirement specifies the exact functionality needed to achieve that impact; documenting requirements in an SRS prevents costly errors later in the project.

\subsection{Questions}
\begin{enumerate}
    \item \textbf{Recall:} According to the IEEE, how is a \emph{software requirement} defined?  Provide the definition in your own words. \newpage
    
    \item \textbf{Understand 1:} In your own description, contrast a \emph{business goal} with a \emph{software requirement}. Highlight what each represents and how they relate to one another. \newpage
    
    \item \textbf{Understand 2:} Explain why a wellcrafted Software Requirements Specification (SRS) helps keep development cost and effort low. Use the ideas from the theory section to support your answer. \newpage
    
    \item \textbf{Apply 1:} The business goal is ``reduce customer support call time by 20 %.'' Write \textbf{one} concrete software requirement that could support this goal. The requirement should be clear, testable, and follow the style of the worked example (e.g., ``The system shall ...''). \newpage
    
    \item \textbf{Apply 2 (numerical):} Assume the current average supportcall duration is \textbf{5 minutes}.
        \begin{enumerate}
            \item Calculate the target average call duration that achieves a 20 % reduction.
            \item Based on that target, write a concrete software requirement that would help the development team meet the goal. The requirement should refer to the target metric you just calculated.
        \end{enumerate} \newpage
    
    \item \textbf{Analyse 1:} Discuss why missing, incomplete, or ambiguous requirements identified early in a project can cause cost overruns later on. Refer to the concepts of hidden requirements, error propagation, and the role of the SRS. \newpage
    
    \item \textbf{Evaluate 1:} Evaluate the statement: ``Investing more effort in a highquality SRS always reduces overall project risk.''* In your evaluation, consider both the benefits described in the theory and any possible limitations or situations where the statement might not hold true. Provide a reasoned judgment. \newpage
\end{enumerate}

\newpage
% ============================================================
% SECTION  2. The Requirements Process
% ============================================================

% !!!  REVIEW FLAG: 2. The Requirements Process  !!!
% [knowledge_assumption] Add a short preface or reference in the prerequisites that introduces regulatory concepts relevant to safetycritical domains, or incorporate a brief explanation into the question itself so that students are not expected to bring external knowledge.
% !!! END REVIEW FLAGS !!!
\section{2. The Requirements Process}
\subsection{Before We Begin}
Structured analysis techniques help us picture how data moves through a system and how information is stored. A dataflow diagram shows the sources, destinations, and transformations of data. An entityrelationship model captures the main objects of interest and the relationships between them. These visual tools let analysts discuss ideas with stakeholders before any code is written. Understanding these basics will let us see why the requirements process can use both formal diagrams and informal conversation.

\subsection{Motivation}
Planning a road trip: you sketch a route, discover road closures, adjust stops, and keep refining the plan as new information appears.  
In software projects, the ``requirements process'' exists to turn vague user wishes into a clear, testable Software Requirements Specification (SRS). It solves the problem of missing or misunderstood needs by continuously checking and refining what the system must do, much as a traveler revises a map when new obstacles appear.

\subsection{Intuition}
Think of the requirements process as a collaborative dialogue. First, we listen to users and gather everything they need the system to accomplish. Next, we write down those needs in a single, organized document. Finally, we review the document with users to catch mistakes early. Because new insights can appear at any stage, the three activities overlap and may repeat, rather than following a straight line.

\subsection{Formal Definition}
The requirements process consists of three core activities:

\begin{enumerate}
\item \textbf{Problem or Requirement Analysis (Elicitation and Analysis)} -- The analyst gathers information about the application's goals, user needs, and any constraints by interviewing users, examining existing manuals or systems, brainstorming, and organizing the collected material. Completeness of the collected information is verified at this stage.
\item \textbf{Requirement Specification} -- All gathered knowledge is documented in the Software Requirements Specification (SRS), forming a single reference for later design, testing, and contractual work.
\item \textbf{Validation} -- The documented requirements are examined with stakeholders to detect errors, ambiguities, or omissions before extensive development proceeds.
\end{enumerate}

The process is \emph{iterative} because validation often uncovers gaps that send the analyst back to further elicitation or analysis. It is \emph{parallel} because analysis, specification, and validation can overlap; for example, a draft requirement may be written while the interview continues, allowing immediate feedback.

Two families of analysis techniques are used:

\begin{itemize}
\item \textbf{Structured approaches} -- These employ formal models such as dataflow diagrams, entityrelationship diagrams, objectoriented analysis, use cases, user stories, or formal mathematical modeling. They provide precise, visual representations that help ensure completeness and consistency.
\item \textbf{Informal approaches} -- These rely on sketches, adhoc diagrams, the analyst's experience, and repeated client reviews. They are flexible and fast, especially when combined with usecase narratives, but they may miss subtle relationships that formal models capture.
\end{itemize}

\subsection{Worked Example}
\textbf{Problem statement:} Conduct an elicitation interview for a small librarymanagement system and derive three highlevel requirements.

\textbf{Step 1 -- Prepare interview questions.}  
\begin{itemize}
\item ``How do users check out books?''  
\item ``What different types of users interact with the system?''  
\item ``What reports do librarians need to generate?''
\end{itemize}

\textbf{Step 2 -- Conduct interview with the librarian.}  
\begin{itemize}
\item The librarian explains that a patron selects a book, presents a library card, and the system records the due date.  
\item She mentions two user types: \emph{Patrons} who borrow books and \emph{Staff} who manage inventory.  
\item She states that weekly reports on overdue books and monthly statistics on mostborrowed titles are required.
\end{itemize}

\textbf{Step 3 -- Organize the answers.}  
\begin{enumerate}
\item \emph{Book checkout process} -- capture patron ID, book ID, and compute due date.  
\item \emph{User role management} -- differentiate between patrons and staff, granting appropriate permissions.  
\item \emph{Reporting functionality} -- generate overduebooks and populartitles reports on specified schedules.
\end{enumerate}

\textbf{Step 4 -- Write highlevel requirements.}  
\begin{enumerate}
\item The system shall allow a patron to check out a book by scanning the patron's card and the book's barcode, automatically recording the checkout date and calculating a due date.  
\item The system shall support two user roles, \emph{Patron} and \emph{Staff}, each with rolespecific access rights.  
\item The system shall produce a weekly overduebooks report and a monthly mostborrowedtitles report for staff review.
\end{enumerate}

\textbf{Final answer:} The interview yields three concise, highlevel requirements covering checkout, role management, and reporting.  

\textbf{What this shows us:} A short, focused interview can quickly capture the essential functionality of a system, and the gathered statements can be directly transformed into clear, verifiable requirements for the SRS.

\subsection{Key Takeaway}
Remember that the requirements process consists of analysis, specification, and validation, and it is deliberately iterative and parallel to catch problems early, using either structured or informal techniques as the situation demands.

\subsection*{Questions}
\textbf{Question 1 (Recall)}  
List and briefly define the three core activities that make up the requirements process as described in the theory section.

\newpage

\textbf{Question 2 (Understanding)}  
In your own words, explain what is meant by the requirements process being \emph{iterative} and \emph{parallel}. Give an example (using any of the activities) that illustrates each characteristic.

\newpage

\textbf{Question 3 (Understanding)}  
Why is it advantageous for analysis, specification, and validation to overlap rather than occur strictly one after another? Discuss how this overlap helps uncover errors early in a project.

\newpage

\textbf{Question 4 (Application)}  
Consider a small online store that sells handmade crafts.  
\begin{enumerate}
\item Formulate a single elicitation question you would ask the store owner to discover a functional need.  
\item Based on a plausible answer, write a concise highlevel requirement that could be added to the SRS.
\end{enumerate}

\newpage

\textbf{Question 5 (Application)}  
Using the same ecommerce scenario, propose an elicitation question aimed at uncovering a nonfunctional requirement (e.g., performance, security, usability).  
Provide the likely stakeholder response and translate it into a clear nonfunctional requirement statement.

\newpage

\textbf{Question 6 (Analysis)}  
Create a sidebyside comparison table that lists at least three advantages and three disadvantages of:
\begin{itemize}
\item Employing a formal dataflow diagram during requirements analysis.  
\item Relying solely on informal interviews without any diagrammatic models.
\end{itemize}
Explain how each approach might affect completeness, stakeholder communication, and project risk.

\newpage

\textbf{Question 7 (Evaluation)}  
A team is developing a safetycritical medical device (e.g., an insulin pump).  
Critically evaluate whether using only informal analysis techniques (e.g., interviews, adhoc sketches) could produce an adequate requirements specification for this domain.  
In your assessment, consider regulatory compliance, risk of omitted requirements, verification needs, and any circumstances where informal methods might still play a role. Provide a reasoned conclusion and suggest a balanced approach.

\newpage
% ============================================================
% SECTION  3. Software Requirements Specification (SRS) Document
% ============================================================

% !!!  REVIEW FLAG: 3. Software Requirements Specification (SRS) Document  !!!
% [undefined_term] Define "API level" in a brief note or footnote, explaining that it refers to the Android SDK version number used by the application.
% !!! END REVIEW FLAGS !!!
\section{3. Software Requirements Specification (SRS) Document}

\subsection{Before We Begin}
Functional requirements describe what the system must do, whereas nonfunctional requirements describe how well it must do it. A measurable performance criterion specifies a numeric target, such as a maximum response time, that can be tested. Design constraints limit choices such as hardware platforms or programming languages. External interfaces detail how users or other systems will interact with the software. Understanding these distinctions prepares us to recognize the four main parts of a complete SRS.  

\subsection{Motivation}
An architect's blueprint: it shows every room (functionality), material limits (design constraints), loadbearing specs (performance), and how utilities connect (external interfaces). Just as a blueprint guides the construction of a safe, usable building, an SRS guides the construction of software that meets all stakeholder needs. Without a clear, structured specification, developers may build the wrong features, miss crucial quality goals, or waste effort on unsupported platforms.  

\subsection{Intuition}
Think of the SRS as a detailed contract between the people who want the software and the people who will build it. It lists what the system shall do (functionality), how fast or reliable it must be (performance), any limits on technology or hardware (design constraints), and how anyone will talk to it (external interfaces). Each item is written so that anyone can check, by observation or testing, whether the software complies. This makes the later stages of design, coding, and testing far smoother and cheaper.  

\subsection{Formal Definition}
The Software Requirements Specification (SRS) is a document that must contain the following four component categories:

\begin{enumerate}
  \item \textbf{Functionality}: a description of all inputs the system accepts, the outputs it produces, the operations it performs, and the behavior it should exhibit when given invalid inputs.  This covers every user category (e.g., administrators, regular users).
  \item \textbf{Performance Requirements}: specifications of dynamic constraints such as response time or throughput, and static constraints such as capacity limits.  Each performance statement must be \emph{measurable}, for example ``response time $\le 2\,$s for $90\%$ of transactions,'' so that testing can confirm compliance.
  \item \textbf{Design Constraints}: explicit restrictions on the technology stack, hardware platforms, applicable standards, reliability targets, and security measures.  These constraints are not optional; they guide architects and developers in selecting solutions that fit the project's context.
  \item \textbf{External Interfaces}: definitions of how users interact with the system (userinterface prototypes, navigation flows) and how the system communicates with other hardware or software components (APIs, data exchange formats).  The user interface is the most visible part and therefore receives special emphasis.
\end{enumerate}

A requirement is \emph{verifiable} when there exists an objective test, inspection, or analysis that can demonstrate whether the requirement has been satisfied.  Vague statements such as ``the system should be fast'' are not verifiable because they lack a measurable target.  Ensuring verifiability prevents the misconception that any wellworded statement is acceptable.  Likewise, the inclusion of nonfunctional items (performance, design, interfaces) is essential; omitting them does not make the SRS complete.  Finally, the SRS is a living document---updates are required whenever stakeholders' needs evolve or new constraints emerge.  

\subsection{Worked Example}
\textbf{Problem statement:}\\
Create a short excerpt of an SRS for a simple todo list mobile application.  The excerpt must contain (a) one functional requirement, (b) one measurable performance requirement, (c) one design constraint, and (d) one description of an external interface.

\textbf{Step 1 -- Identify a functional requirement.}\\
The application must allow a user to add a new task to their personal todo list.

\textbf{Step 2 -- Write the functional requirement.}\\
\textit{FR1:} The system shall provide a screen where the user can enter a task description and press a ``Create'' button; upon activation, the task shall be stored and immediately appear in the user's task list.

\textbf{Step 3 -- Identify a performance requirement.}\\
We need a responsetime target for adding a task.

\textbf{Step 4 -- Write the performance requirement.}\\
\textit{PR1:} The system shall complete the ``Create Task'' operation within $1\,$s for $95\%$ of user actions measured over a 30day period.

\textbf{Step 5 -- Identify a design constraint.}\\
The app must run on modern Android devices.

\textbf{Step 6 -- Write the design constraint.}\\
\textit{DC1:} The application shall be compatible with Android version $10$ (API level $29$) or higher.

\textbf{Step 7 -- Identify an external interface.}\\
The app will expose a RESTful service for creating tasks from other clients.

\textbf{Step 8 -- Write the external interface description.}\\
\textit{EI1:} The system shall provide a POST endpoint at \texttt{/api/tasks} that accepts a JSON payload \{``description'': string\}; a successful request shall return HTTP status $201$ with the created task's identifier.

\textbf{Final answer:}\\
\begin{itemize}
  \item \textbf{FR1:} The system shall provide a screen where the user can enter a task description and press a ``Create'' button; upon activation, the task shall be stored and immediately appear in the user's task list.
  \item \textbf{PR1:} The system shall complete the ``Create Task'' operation within $1\,$s for $95\%$ of user actions measured over a 30day period.
  \item \textbf{DC1:} The application shall be compatible with Android version $10$ (API level $29$) or higher.
  \item \textbf{EI1:} The system shall provide a POST endpoint at \texttt{/api/tasks} that accepts a JSON payload \{``description'': string\}; a successful request shall return HTTP status $201$ with the created task's identifier.
\end{itemize}

\textbf{What this shows us:}  
A concise SRS excerpt can capture all four required component categories---functional behavior, a quantifiable performance target, a concrete technology limitation, and a clear external interaction---demonstrating how each piece contributes to a complete, verifiable specification.  

\subsection{Key Takeaway}
An SRS must contain clear, measurable statements for functionality, performance, design constraints, and external interfaces; only such a complete and verifiable specification can guide successful software development.

\subsection{Questions}

\noindent\textbf{Question 1:}  
List and briefly name the four major component categories that must be included in a complete Software Requirements Specification (SRS). \newpage

\noindent\textbf{Question 2:}  
In your own words, explain why every requirement written in an SRS should be \emph{verifiable}. \newpage

\noindent\textbf{Question 3:}  
Give one concrete example for each of the four SRS component types (Functionality, Performance Requirements, Design Constraints, External Interfaces) that shows how a verifiable statement differs from a vague, nonverifiable one. \newpage

\noindent\textbf{Question 4:}  
Write a \textbf{performance requirement} for a mobile chat application that specifies the messagedelivery latency in a measurable way. Include a numeric target and the proportion of messages to which it applies. \newpage

\noindent\textbf{Question 5:}  
The same chat application must handle at least 10 000 concurrent users during peak hours. Write a \textbf{static performance requirement} that states this capacity using a measurable figure. \newpage

\noindent\textbf{Question 6:}  
Adapt the functional requirement from the worked example (FR1) to the chat application: the system shall allow a user to send a text message. Write this functional requirement in the same style as the example, including the expected behavior for invalid input (e.g., empty message). \newpage

\noindent\textbf{Question 7:}  
Analyse the possible consequences of \emph{omitting design constraints} from an SRS. Discuss at least two ways this omission could increase later development effort, using examples such as technologystack selection or security compliance. \newpage

\noindent\textbf{Question 8:}  
Critically evaluate the claim: ``A highly detailed SRS always reduces the chance of project failure.''* In your answer, consider both the potential benefits of detail and the risks of overspecification, drawing on the concepts of correctness, completeness, and verifiability presented in the section. \newpage

\newpage
% ============================================================
% SECTION  4. Specifying Functional Requirements (Use Cases & User Stories)
% ============================================================

\section{4. Specifying Functional Requirements (Use Cases \& User Stories)}

\subsection{Before We Begin}
Before we begin, remember that an \emph{actor} is any person or external system that interacts with the software we are describing. In agile methods, a \emph{user story} is a short, structured sentence that captures who wants what and why. Understanding these ideas lets us create clear, userfocused specifications. With that groundwork, we can explore how to describe what a system must do for its users.

\subsection{Motivation}
A theater script: actors (users) follow a series of scenes (use cases) to tell a story, while a quick backstage note (user story) tells the director what a performer needs. Defining functional requirements tells developers exactly what the system should do, preventing costly misunderstandings later. Use cases give a stepbystep picture of interactions, while user stories provide a lightweight way to capture the same intent from a user's perspective.

\subsection{Intuition}
Think of a use case as a short play: a recognizable character (the actor) steps on stage, pursues a goal, and follows a main series of actions that lead to success, possibly with alternative twists. A user story, by contrast, is like a sticky note left for the director: it names the performer, states the desired action, and explains the benefit, without spelling out every line of dialogue. Both help the team agree on what the system must deliver, but at different levels of detail.

\subsection{Formal Definition}
A \textbf{functional requirement} specifies a behavior that the system must provide to its users. One common way to capture functional requirements is a \textbf{use case}. A use case is a documented scenario that includes the following four essential elements:

\begin{enumerate}
  \item \textbf{Name/Goal}: a unique label that identifies the purpose of the scenario.
  \item \textbf{Actor}: the external person or system that initiates the interaction.
  \item \textbf{Entry condition}: the state of the system that must be true before the scenario can start.
  \item \textbf{Interaction sequence} (or main flow): the ordered steps that the actor and system perform to achieve the goal, together with any alternative or exceptional flows.
\end{enumerate}

Use cases may be linked together. An \emph{Include} relationship means that one use case invokes another as a mandatory step; the included use case is always executed as part of the primary scenario. An \emph{Extend} relationship adds optional behavior that occurs only under certain conditions, modifying or expanding the primary flow.

A \textbf{user story} is a lightweight agile artifact that captures the same information using the template ``As a \emph{type of user}, I want \emph{some action} so that \emph{some benefit}.''
It does not enumerate stepbystep interactions or alternative flows; instead, it records who needs what and why, leaving details for later refinement. Unlike a UI mockup, a use case describes behavior, not visual layout, and a user story does not contain full acceptance criteria---it simply states the desired outcome.

\subsection{Worked Example}
\textbf{Problem statement:} Create a complete use case for an online bookstore's ``Place Order'' functionality (actor: customer, goal: purchase items, entry condition: cart not empty, main flow: review cart, provide payment, confirmation). Then write an accompanying user story: \emph{As a customer, I want to receive an email receipt after purchase so that I have proof of transaction.}

\textbf{Step 1 -- Identify the four elements of the use case.}
\begin{itemize}
  \item \textbf{Name/Goal}: Place Order -- Customer purchases selected items.
  \item \textbf{Actor}: Customer.
  \item \textbf{Entry condition}: The shopping cart contains at least one item.
  \item \textbf{Interaction sequence (main flow)}:
\end{itemize}

\textbf{Step 2 -- Write the main flow.}
\begin{enumerate}
  \item Customer selects ``Proceed to Checkout'' from the cart page.
  \item System displays the order review page with a summary of items, quantities, prices, and total cost.
  \item Customer confirms the items and clicks ``Continue to Payment.''
  \item System prompts for payment information (credit card or alternative).
  \item Customer enters payment details and submits the form.
  \item System validates payment, processes the transaction, and generates an order number.
  \item System shows an order confirmation page with the order number and estimated delivery date.
\end{enumerate}

\textbf{Step 3 -- Add alternative/exception flows (briefly).}
\begin{itemize}
  \item If payment validation fails, system displays an error message and returns to the payment entry step.
  \item If the customer aborts at any step, the system cancels the order and returns to the cart.
\end{itemize}

\textbf{Step 4 -- Write the user story.}
\begin{quote}
As a \textbf{customer}, I want to receive an \textbf{email receipt} after purchase so that I have \textbf{proof of transaction}.
\end{quote}

\textbf{Final answer:} The use case ``Place Order'' captures the complete endtoend purchasing process, while the user story records a specific postpurchase need (email receipt) in the concise ``AsIWantSoThat'' format.

\textbf{What this shows us:} Use cases give a detailed, stepbystep view of system behavior, whereas user stories succinctly express a user's desire and its benefit, both serving complementary roles in functional specification.

\subsection{Key Takeaway}
Use cases describe detailed interaction flows with actors, goals, entry conditions, and main steps, while user stories provide a brief ``whowantswhatwhy'' snapshot of the same requirement.

\subsection*{Questions}
\begin{enumerate}
  \item \textbf{(Recall)} List the four essential elements that must be included in a usecase description according to the theory section. 
  \newpage
  \item \textbf{(Understand)} In your own words, explain the difference between an \emph{Include} relationship and an \emph{Extend} relationship in use cases. Mention when each is typically used. 
  \newpage
  \item \textbf{(Understand)} Using the ``Place Order'' example from the worked example, give one concrete instance of an \emph{Include} relationship and one concrete instance of an \emph{Extend} relationship. Briefly describe the purpose of each relationship in that context. 
  \newpage
  \item \textbf{(Apply)} Write a user story for a fitnesstracking app that captures a user's desire to view their weekly progress. Use the template ``As a \emph{type of user}, I want \emph{some action} so that \emph{some benefit}.''
  \newpage
  \item \textbf{(Apply)} From the user story you created in the previous question, identify the three components (type of user, action, benefit). Count the number of words in each component and give your answer in the form  
  \[
  \text{User: } \_ \text{ words},\; \text{Action: } \_ \text{ words},\; \text{Benefit: } \_ \text{ words}.
  \] 
  \newpage
  \item \textbf{(Apply)} Imagine you need to add an ``Set Weekly Goal'' use case for the same fitnesstracking app. Provide a brief interaction sequence (main flow) consisting of exactly three steps that would allow the user to set and confirm a weekly activity goal. 
  \newpage
  \item \textbf{(Analyse)} Analyse the advantages and disadvantages of relying solely on user stories versus employing detailed use cases for a mediumsize web project. Consider factors such as communication, clarity, effort, and maintenance. 
  \newpage
  \item \textbf{(Evaluate)} Evaluate whether a use case that omits an explicit entry condition can be acceptable in a Software Requirements Specification (SRS). Support your answer with arguments about completeness, traceability, and potential risks. 
  \newpage
\end{enumerate}

\newpage
% ============================================================
% SECTION  5. Minimum Viable Product (MVP) & Prioritization
% ============================================================

\section{5. Minimum Viable Product (MVP) \& Prioritization}
\subsection{Before We Begin}
Lean Startup thinking encourages building only what is needed to test a business hypothesis as quickly as possible. Early validation means releasing a small, usable version of a product, observing real user behaviour, and deciding whether to continue, pivot, or stop. Without this approach a team may invest heavily in features that users never want. Understanding that the goal is rapid learning, not delivering a polished final product, prevents the mistaken belief that an MVP is simply a lowquality version. This idea leads directly to how we prioritize requirements for an MVP.  

\section{Motivation}
Building a sandcastle: you start with just enough walls to hold water before adding decorative towers. In software development the same principle applies: we create a product that contains only the features essential to prove that the market problem exists and that users find value. By limiting scope we avoid ``requirements bloating'' and ``requirements creep'' that would waste time and money on untested ideas. The resulting Minimum Viable Product (MVP) lets us gather real feedback fast and steer future development accordingly.  

\subsection{Intuition}
Imagine you have a long wishlist of things a new app could do. Before writing any code you sort that list into three buckets: features that are absolutely necessary to solve the core problem (Essential), features that would make the product nicer but are not required for the first test (Desirable), and features that can wait until the product is proven and the team has resources (Deferable). An MVP is simply the software that contains only the Essential bucket; it must still work reliably, look decent enough for users, and be built on an architecture that can later accommodate the other buckets.  

\subsection{Formal Definition}
A \textbf{requirement priority level} classifies each functional requirement according to its importance for the first market test:

\begin{itemize}
  \item \textbf{Essential} (also called ``Musthave''): a requirement without which the product cannot fulfill its primary purpose or provide any meaningful value to early users.
  \item \textbf{Desirable}: a requirement that improves the user experience or adds competitive advantage but is not required to validate the core hypothesis.
  \item \textbf{Deferable}: a requirement that is useful or nicetohave but can be postponed until after the MVP has been released and validated.
\end{itemize}

A \textbf{Minimum Viable Product (MVP)} is a software system that:

\begin{enumerate}
  \item implements \emph{all} Essential requirements and \emph{no} requirements from the Desirable or Deferable categories,
  \item delivers \emph{value} to early adopters (the product solves the problem it set out to address),
  \item is \emph{reliable} enough that users can complete the essential tasks without crashes or data loss,
  \item presents a \emph{good enough} user interface so that users can understand and use the essential features,
  \item is built on an \emph{extensible architecture} that permits adding Desirable and Deferable features in later releases.
\end{enumerate}

These statements preempt common misconceptions: an MVP is not a halffinished, lowquality product; it is a deliberately limited but fully functional slice of the final vision. Including highvalue features that are not Essential defeats the purpose of rapid validation, and prioritization is not a onetime planning step---it directly determines what is built now and what can be safely postponed.  

\subsection{Worked Example}
\textbf{Problem statement:} Design the scope of an MVP for a small ridesharing application.

\textbf{Step 1 -- List candidate requirements.}
\begin{itemize}
  \item User registration (create an account).  
  \item Request a ride (passenger enters pickup and dropoff locations).  
  \item Driver acceptance (driver receives request and can accept).  
  \item Ridesharing options (passenger can request to share a ride with others).  
  \item Inapp payment (passenger pays through the app).
\end{itemize}

\textbf{Step 2 -- Categorize each requirement.}
\begin{itemize}
  \item \textbf{Essential:} user registration, request a ride, driver acceptance.  
  \item \textbf{Desirable:} ridesharing options.  
  \item \textbf{Deferable:} inapp payment.
\end{itemize}

\textbf{Step 3 -- Define the MVP scope.}  
The MVP will implement exactly the three Essential requirements: a passenger can sign up, request a ride, and a driver can accept that request. No ridesharing or payment features are included at this stage.

\textbf{Step 4 -- Explain how the MVP validates the market hypothesis.}
The hypothesis is ``Users need a quick, reliable way to match with a driver for a single ride.'' By delivering only registration, ride request, and driver acceptance, the product can be released to a small group of early adopters. Their usage patterns (how many rides are requested, acceptance rates, repeat usage) provide concrete evidence whether the core matching service is valuable. If the data is positive, the team can safely add the Desirable ridesharing option and later the Deferable payment feature.

\textbf{Final answer:} The MVP consists of the three Essential features listed above, and it serves to test the fundamental matching hypothesis before any extra functionality is built.

\textbf{What this shows us:} Prioritizing requirements lets us build the smallest possible product that still answers the key business question, thereby avoiding unnecessary work and enabling fast learning.  

\subsection{Key Takeaway}
An MVP includes \textbf{all} Essential requirements and nothing else, providing a reliable, valuable product that lets a team test its core market hypothesis quickly.  

\subsection*{Questions}
\newpage
\textbf{1. (Recall)} According to the formal definition of requirement priority levels in the MVP discussion, \emph{list the three priority categories} and give a brief description of each (Essential, Desirable, Deferable).  

\newpage
\textbf{2. (Understand)} In your own words, \emph{explain what ``requirement creep'' means} and why allowing new requirements to be added after the MVP scope has been set can weaken the purpose of an MVP.  

\newpage
\textbf{3. (Understand)} Describe how \emph{requirement creep} can affect the speed of learning from an MVP release. Include in your answer the impact on feedback cycles and the risk of building features that have not been validated.  

\newpage
\textbf{4. (Apply)} You are designing a simple notetaking app.  
Identify \emph{two Essential requirements} that must be included in the MVP and \emph{two Desirable requirements} that could be deferred to later releases. Briefly justify why each requirement belongs to its chosen category.  

\newpage
\textbf{5. (Apply)} The following five candidate requirements are proposed for a notetaking app:

\begin{enumerate}
  \item Create, edit, and delete notes.  
  \item Sync notes across multiple devices.  
  \item Search notes by keyword.  
  \item Export notes to PDF.  
  \item Darkmode UI theme.
\end{enumerate}

Classify each requirement as \textbf{Essential}, \textbf{Desirable}, or \textbf{Deferable} for the MVP. Then calculate the \emph{percentage of requirements that are Essential} in this list.  

\newpage
\textbf{6. (Apply)} Assume an MVP team decides to implement exactly three Essential features out of a total of eight identified requirements.  
How many features will be \emph{omitted} (i.e., not implemented) in the MVP? Show your calculation.  

\newpage
\textbf{7. (Analyse)} Analyse the risks associated with launching an MVP that \emph{does not provide a reliable user interface}. Consider aspects such as user trust, data collection, and the ability to obtain meaningful feedback.  

\newpage
\textbf{8. (Evaluate)} Evaluate whether an MVP that \emph{omits datalogging capabilities} (e.g., no recording of user actions, usage statistics, or error reports) can still satisfy the viability criteria defined for an MVP. Discuss the tradeoffs and under what conditions the omission might be acceptable or detrimental.  

\newpage

\newpage
% ============================================================
% SECTION  6. Project SRS Development Process & Document Structure
% ============================================================

% !!!  REVIEW FLAG: 6. Project SRS Development Process & Document Structure  !!!
% [knowledge_assumption] Add a brief description of the role of traceability in testing activities to the prior sections so that students can connect traceability to test case creation, requirement verification, and defect tracking.
% !!! END REVIEW FLAGS !!!
\section{6. Project SRS Development Process \& Document Structure}

\subsection{Before We Begin}
A Large Language Model (LLM) is a computer program that has been trained on vast amounts of text so it can generate new sentences that sound like a human wrote them. When you ask an LLM to write a draft, it quickly produces a first version of the document based on the information you give it. This first version usually needs a human to check it, add missing details, and correct any mistakes. Understanding that an LLM can help you start a document, but not finish it, prepares us for its role in the SRS drafting stage.  

\section{Motivation}
Cooking a multicourse meal: you first discuss the menu with the guest (sponsor), then write a draft recipe list (draft), and finally tastetest each dish together before serving (final review). Just as a wellplanned meal prevents disappointment at the table, a structured threestep process for a Software Requirements Specification (SRS) prevents costly misunderstandings later in a project. It gives the sponsor a clear view of what will be built, helps the development team capture every needed feature, and ensures that nothing is missed before any coding begins.  

\subsection{Intuition}
Think of an SRS as a detailed recipe for a software product. First you sit down with the person who will eat the meal (the sponsor) to learn what flavors they want and any dietary restrictions. Next you write down a rough list of ingredients and steps, perhaps getting a quick draft from an assistant (the LLM). Finally, you go through the recipe line by line with the guest, tasting and adjusting each part until everyone agrees the dish is ready to be cooked.  

\subsection{Formal Definition}
The Project SRS Development Process consists of three sequential activities:

\begin{enumerate}
    \item \textbf{Analyze}: Meet the sponsor to capture the business goal, the types of users who will interact with the system, and the overall scope of the project.
    \item \textbf{Draft}: Produce an initial SRS document. An LLM may be employed to generate a firstpass outline of features, after which the analyst refines the text, adds missing requirements, and corrects any inaccuracies.
    \item \textbf{Finalize}: Perform a meticulous linebyline review together with the sponsor, confirming that every stated requirement matches the sponsor's intent before the document is declared complete.
\end{enumerate}

A standard SRS for a course project contains the following toplevel sections:

\begin{itemize}
    \item \textbf{Title of the project}
    \item \textbf{Sponsor details} -- organization name and contact information
    \item \textbf{Development Team details}
    \item \textbf{Introduction} -- purpose, system objectives, user types, and constraints
    \item \textbf{Functional requirements overview} -- a usecase table for each user category, with descriptions and priority rankings
    \item \textbf{NonFunctional requirements} -- userinterface expectations, performance criteria, and design/technology stack
    \item \textbf{Further Details} -- expanded descriptions of each use case
\end{itemize}

Note that the draft is not the final product, sponsor involvement is needed throughout, and every listed section plays a role in creating a complete, testable specification.  

\subsection{Worked Example}
\textbf{Problem statement:} Produce a skeleton SRS for a campus eventregistration system.

\textbf{Step 1 -- Title}\\
\verb|Title: Campus Event Registration System|

\textbf{Step 2 -- Sponsor}\\
\verb|Sponsor: University Student Activities Office|\\
\verb|Contact: alice.smith@university.edu|

\textbf{Step 3 -- Development Team}\\
\verb|Team Lead: Bob Lee|\\
\verb|Developers: Charlie Kim, Dana Patel|

\textbf{Step 4 -- Introduction}\\
\begin{itemize}
    \item Purpose: Enable students to browse, register for, and receive confirmation of campus events.
    \item System Objectives: Provide a searchable event catalog, secure registration, and automatic email notifications.
    \item User Types: Student (regular user), Event Organizer (admin), Guest Lecturer (external user).
    \item Constraints: Must integrate with existing university single signon (SSO) system.
\end{itemize}

\textbf{Step 5 -- Functional Requirements Overview (UseCase Table)}  

\begin{tabular}{|l|l|l|l|}
\hline
\textbf{Use Case ID} & \textbf{Actor} & \textbf{Goal} & \textbf{Priority} \\
\hline
UC001 & Student & Register for an event & High \\
\hline
UC002 & Event Organizer & Create a new event listing & Medium \\
\hline
UC003 & Guest Lecturer & View event details without registration & Low \\
\hline
\end{tabular}

\textbf{Step 6 -- NonFunctional Requirements}\\
\begin{itemize}
    \item UI: Responsive web design for desktop and mobile browsers.
    \item Performance: Page load time $< 2$ seconds under normal load.
    \item Design/Tech Stack: Frontend React, backend Node.js, PostgreSQL database.
\end{itemize}

\textbf{Step 7 -- Further Details}\\
\begin{itemize}
    \item Detailed flow for UC001: login $\rightarrow$ search events $\rightarrow$ select event $\rightarrow$ click ``Register'' $\rightarrow$ confirmation email.
    \item Detailed flow for UC002: login as organizer $\rightarrow$ access admin dashboard $\rightarrow$ fill event form $\rightarrow$ submit $\rightarrow$ event appears in catalog.
\end{itemize}

\textbf{Final answer:} The skeleton above contains every required toplevel heading and a representative entry in the functional requirements table.  

\textit{What this shows us:} By following the threestep process and using the standard SRS structure, even a new analyst can quickly produce a clear, organized specification that the sponsor can review line by line.  

\subsection{Key Takeaway}
The threestep \textit{analyzedraftfinalize} process, together with the prescribed SRS sections, ensures that a software specification is complete, verified by the sponsor, and ready for implementation.

\subsection*{Questions}
\textbf{Question 1 (Recall):}\\
List the three main steps in the Project SRS Development Process as defined in the theory section.\newpage

\textbf{Question 2 (Understand):}\\
In your own words, describe why conducting a linebyline review of the SRS together with the sponsor is essential before the document is declared final.\newpage

\textbf{Question 3 (Understand):}\\
How does the linebyline review help ensure that the SRS meets the sponsor's business goal, usertype expectations, and scope constraints?\newpage

\textbf{Question 4 (Apply):}\\
Create an outline for the SRS of a simple budgeting app. List all required toplevel sections (Title, Sponsor details, Development Team, Introduction, Functional Requirements Overview, NonFunctional Requirements, Further Details) and provide brief bulletpoint examples for at least two of those headings (e.g., two functional usecases with descriptions and priority rankings).\newpage

\textbf{Question 5 (Apply):}\\
For the same budgeting app, specify two nonfunctional requirements that include numerical values. For example, give a performance target (e.g., pageload time) and a reliability metric (e.g., uptime percentage). Write them in the format used in the worked example.\newpage

\textbf{Question 6 (Analyse):}\\
Analyse the potential drawbacks of relying solely on an LLMgenerated draft of the SRS without any analyst verification. Consider issues such as missing requirements, ambiguity, and alignment with sponsor intent.\newpage

\textbf{Question 7 (Evaluate):}\\
Evaluate whether omitting a ``Traceability'' subsection from an SRS could cause problems during later testing phases. Discuss the possible impacts on test case creation, requirement verification, and defect tracking.\newpage

\end{document}
